\documentclass[11pt,letterpaper]{exam}
\usepackage[top=1in,bottom=1in,left=1in,right=1in]{geometry}
\usepackage{amsmath,amssymb,graphicx}
\usepackage{xcolor,hyperref,enumitem,booktabs}
\newcommand{\blue}[1]{\textcolor{blue}{#1}}

\begin{document}

\begin{center}
\textbf{\Large ECON 2181 – International Trade Theory \\[4pt]
Data Lab Project: Comparative Statics in the Specific Factors Model}
\end{center}

\bigskip

\noindent
The world economy has two countries $i\in\{H,F\}$ and two goods $g\in\{M,A\}$.  
Notation follows the notebook exactly:
\[
\begin{aligned}
Y_{i,M} &= Z_{i,M}\,K_i^{\beta_i}\,L_{i,M}^{\,1-\beta_i}, &
Y_{i,A} &= Z_{i,A}\,T_i^{\beta_i}\,L_{i,A}^{\,1-\beta_i}, \\
L_{i,M}+L_{i,A}&=\bar L_i, &
p &\equiv P_M/P_A.
\end{aligned}
\]
The world‐equilibrium relative price $p^*$ solves $RS(p)=RD(p)$.

\bigskip\hrule\bigskip

\begin{questions}
\question Run the notebook with baseline parameters. Record in specific variables:
\[
p^*,\quad w_H,w_F,\quad Y_{i,M},Y_{i,A},\quad L_{i,M},L_{i,A}, 
\]

\question Increase manufacturing productivity by 20\% in both countries:
\[
Z'_{i,M}=1.2\,Z_{i,M}.
\]
\begin{parts}
\part Recompute $p^*$, the wage rates, labor allocations, and outputs. Record them in new variables that will allow you to compare the baseline and new equilibria.
\part Compare $L_{i,M}, L_{i,A}$ before and after the shock. You do not need to plot charts. Is this consistent with what we saw in class?
\part Plot relative world production $( Y_{H,M} +  Y_{F,M} ) / ( Y_{H,A} +  Y_{F,A} )$ and $p^*$ before and after the change. What happens with relative output $Y_M/Y_A$ and relative prices $P_M/P_A$? Explain intuitively what happens in terms of supply and demand. 
\part Compute real wages $w_i/P_M$ and $w_i/P_A$ and compare to baseline. Is there a change in real wages? Is the consumer better off or worse off?
\end{parts}
\end{questions}
\end{document}
